\label{sec:repro}

\paragraph{Code and frozen states.} Code, manifests, aggregated tables, statistics and figures are in the repository
\texttt{caioribeiro99/doe\_nbi\_hpo\_project}. Two annotated tags freeze the study: \texttt{pco213-\allowbreak postwork-\allowbreak r10}
the ten-replication analysis and \texttt{pco213-\allowbreak postwork-\allowbreak r30} (commit \texttt{b3ed050}) the
thirty-replication analysis reported here (\cref{sec:results:r30} compares the two tags). The manuscript, figures and
tables are on the branch \texttt{paper/\allowbreak surrogate-\allowbreak nbi-\allowbreak ensemble-\allowbreak weighting}, cut from that tag, together with the results computed after it from the stored outputs of the
thirty-replication run, among them the NSGA-II baseline, the anchor-injection control and the blend test of
\cref{tab:edge}. A single checkpointed ten-stage runner drives the benchmark and
resumes an interrupted run without recomputation; its stages, the resumption planner, the aggregation and statistics
scripts and the coverage of the library's 83 unit tests are in supplementary Section~S1. The commit history records
that the library, the benchmark runner and the analysis and asset-building scripts were developed with the assistance
of Anthropic's Claude models (Claude Fable 5, Claude Fable 5.1 and Claude Opus 5) used through Claude Code; the
library's behaviour is checked by those unit tests.

\paragraph{Determinism and environment.} Every stochastic component is seeded as $\text{20260904} + r$ for
replication $r$; the offsets, the seed map, the configuration history and the resolved library versions are in the
master manifest (supplementary Section~S1). The environment is Python 3.11.15 with NumPy 2.4.6 \citep{harris2020numpy}, pandas 3.0.5,
SciPy 1.17.1 \citep{virtanen2020scipy}, scikit-learn 1.9.0 \citep{pedregosa2011scikit}, statsmodels 0.15.0 and
XGBoost 3.2.0 \citep{chen2016xgboost}, on macOS (Apple M4 Max, arm64), with eight worker threads. Execution provenance --- stage time and the fit, design, subproblem and objective-evaluation counts --- is in
supplementary Section~S1. Two stages were retried, both in the first ten replications; none in the extension.

\paragraph{Data availability and redistribution limits.} The raw datasets are not redistributed: the UCI
default-of-credit-card-clients data are downloadable from the UCI Machine Learning Repository, while the three Kaggle
datasets are subject to competition rules and must be obtained from Kaggle by the user. Per-dataset manifests record
the provenance (source, acquisition command, expected invariants and the SHA-256 checksum of the file actually
used), so a copy of the same file can be verified byte-for-byte (supplementary Section~S1). All three Kaggle training
files were obtained from public mirror datasets hosted on Kaggle and validated against those invariants; the
manifests name each mirror. Competition test labels are never used. Porto Seguro was subsampled to 200{,}000 rows
by a stratified draw with the recorded seed. The per-replication raw artifacts total about 829~MB and are
deliberately not version-controlled, but are reconstructible from the seeds and the runner, except for the
wall-clock costs $c_i$, which are re-measured on each run, and the results that depend on them; the measured costs
are committed in \texttt{tables/inference\_costs.csv}.

\paragraph{What a replication requires.} Reproducing the aggregate results from the published tables and statistics
requires only the repository; from raw data it requires the four datasets, the pinned environment and approximately 81
hours of single-workstation compute, or proportionally less for a subset of datasets or replications, since the
runner is resumable and parameterized by dataset and replication count.
